\section{Tasks, Architectures, and Optimization Methods}\label{sec:matched}
An \emph{architecture} specifies the operations and which scalar coefficients may change. A \emph{rectified linear unit} (ReLU) applies $R(u)=\max(0,u)$ to an input $u$. A multilayer perceptron (MLP) applies a sequence of affine maps (weighted sums plus biases) and activations. A recurrent neural network (RNN) carries a state vector from one step to the next and reuses the same weights at each step. Table~\ref{tab:architecture} lists every primary graph. The counts include all trainable biases and shared weights, with each stored scalar counted once.

\begin{table}[htbp]\centering\small
\caption{Architectures in the primary comparison. Sizes and parameter counts correspond left to right. A width counts hidden units or recurrent state coordinates. For sorting, size counts consecutive residual blocks, each acting on the same four wires.}\label{tab:architecture}
\begin{tabular}{@{}llrr@{}}\toprule
Target & Network & Sizes & Trainable parameters\\\midrule
Square, sine, positive self-power, & One hidden ReLU layer & $1,4,16$ & $5,14,50$\\
\quad hyperbolic tangent & with a direct input path & & \\
Addition, subtraction & Two-input ReLU MLP & $1,2,4$ & $5,9,17$\\
Running sum & Linear RNN with input & $1,2$ & $3,8$\\
Sort four numbers & Four-wire residual ReLU network & $4,5,6$ & $32,40,48$\\
Fibonacci numbers & Autonomous linear RNN & $1,2,3$ & $3,8,15$\\
Triangular numbers & Autonomous linear RNN & $2,3,4$ & $8,15,24$\\\bottomrule
\end{tabular}\end{table}

\paragraph{Graph definitions.}
The scalar targets are $x^2$, $\sin x$ in radians, $x^x$ on positive inputs, and the hyperbolic tangent $\tanh x$. For each hidden unit $j$ let $u_j,v_j,c_j$ denote its input slope, offset and output weight, and let $w$ be the width. With a direct linear path of slope $a$ and offset $b$, the scalar prediction is $ax+b+\sum_{j=1}^w c_j R(u_jx+v_j)$. Each unit introduces at most one change in slope, called a hinge. The addition and subtraction networks take two inputs, have no direct path and learn two input weights, a bias and an output weight per hidden unit, plus a final output bias.

For running sum the scalar input at tick $t$ is $x_t$. With a state of width $d$, a learned transition matrix $A$, input column $b$ and readout column $v$, the graph computes $h_t=Ah_{t-1}+b x_t$ from $h_0=0$ and predicts $v^\top h_t$ at every tick. Here $\top$ denotes transpose. For Fibonacci and triangular numbers the learned column $b$ is instead the initial state $h_0$, there are no inputs after initialization, and each tick applies $h_{t+1}=Ah_t$. At integer index $n\ge0$, Fibonacci numbers start at $0,1$ and each later term is the sum of its two predecessors. Triangular numbers equal $n(n+1)/2$. For sorting, the state $h$ initially holds the four inputs. Each block updates it to $h+vR(u^\top h)$ with its own learned four-component columns $u,v$. The final four values should give the original inputs in ascending order.

\paragraph{How we checked nonconvexity.}
For each of the 87 fixed-network settings, we exhibit two weight vectors whose midpoint has loss strictly greater than the average endpoint loss. Each witness uses two trainable coefficients that multiply along an active path. Swapping those coefficients preserves the endpoint prediction while changing the midpoint prediction. Exact rational evaluation on the actual training observations certifies the violation. Appendix~\ref{app:nonconvex-details} supplies the construction and its scope. The additional architecture-selection experiment can choose a different graph and is reported separately.



\paragraph{Initialization and observations.}
We used three random starts per task and architecture. Under Xavier-uniform initialization a layer with $p$ inputs and $q$ outputs draws its weights uniformly between $-\sqrt{6/(p+q)}$ and $\sqrt{6/(p+q)}$. Designated biases start at zero. All primary methods start with the same initial arrays and 12 central labeled inputs for targets with sampled inputs. A running-sum example is a length-six stream and a sorting example is a four-item list. The two integer sequences instead use the nine fixed indices 0 through 8. A new seed changes both the random training data and the weights where inputs are sampled. Test inputs are fixed across sizes and seeds. The three sizes (two for running sum) share a single WSP request, and the three integer-sequence starts share the same training indices.

\paragraph{Objective and methods.}
For observed scalar prediction $\hat y$ with label $y$, the scaled residual is $(\hat y-y)/\max(1,|y|)$. The common training objective is half the mean squared scaled residual, over every output coordinate. Both WSP prompts explicitly prioritize this objective and request extrapolation among similarly good fits. An \emph{endpoint} is the weight array a method returns using only the specified training information. The original four endpoints are the two WSP conditions and two numerical solvers.
\begin{itemize}
\item \textbf{One-shot WSP.} In the \emph{Examples only} condition, Astra sees the graph, initial weights, intended input range and examples. In \emph{Formula and examples}, it also sees the target rule. It returns one array per size in a single reply without tools or scoring feedback.
\item \textbf{Iterative WSP.} We retain both original information conditions. At each active size, Astra receives the selected weights, exact-match counts, errors and up to three worst labeled examples from each of the training and development sets. The development pool contains 256 new central inputs per sampled task and start, or indices 9--16 for the integer sequences. A revision replaces the current weights only when its development loss is strictly lower or, at equal development loss, its training loss is lower. An invalid proposal counts as a rejection. Each size stops when every training and development example is exact, after two consecutive rejections, or after three revision attempts. This is a finite stopping rule, not a guarantee of convergence. Requests for different cases may run concurrently. Each revision uses a fresh Astra session with the same task and selected weights; the prompt forbids tools, and we inspect the recorded events for tool use. Test inputs and scores never enter a revision. Appendix~\ref{app:feedback} describes the audit and controls.
\item \textbf{Gradient descent and gradient descent followed by L-BFGS.} Full-batch gradient descent uses the derivative computed by backpropagation and a training-loss line search for up to 800 updates. From that endpoint, the limited-memory Broyden--Fletcher--Goldfarb--Shanno method (L-BFGS) approximates curvature with ten recent pairs of weight and gradient changes, for up to 1,000 iterations \cite{liu1989limited}. We score the two endpoints separately.
\item \textbf{Adam followed by BFGS.} Adaptive moment estimation (Adam) scales each update using running averages of the gradient and its square \cite{kingma2015adam}. From Xavier we run 800 full-batch steps and then up to 1,000 iterations of the Broyden--Fletcher--Goldfarb--Shanno method (BFGS), which estimates local curvature from changes in gradients. The solver returns the lower training loss of its two stages.
\item \textbf{Trust-region least squares.} Starting at the identical Xavier weights, this solver uses the matrix of first derivatives of individual residuals (the Jacobian) to select local steps. It allows up to 250 residual evaluations followed by up to ten local polishing steps accepted on training loss. 
\end{itemize}
Appendix~\ref{app:limitedmemory} also reports L-BFGS directly from Xavier and from the 800-step Adam endpoint. The primary numerical solvers fit the original training labels. A separate control fits both gradient baselines to the entire labeled training and development pool.

The archive also scores the initial Xavier array. Astra brings pretrained knowledge of familiar formulas and operates under a different computational budget. Examples-only prompts disclose the deployment range and a task-specific graph. A match on visible information therefore controls the experiment more narrowly than a match on prior knowledge or compute.

\paragraph{Evaluation.}
We evaluate root mean squared error (RMSE) in target units over all scalar outputs. An exact prediction equals its reference in the standard 64-bit floating-point format (binary64), treating both signs of zero as equal. For running sums and sorting, one example counts as exact only if every coordinate matches. For the integer sequences we count matching indices on a single trajectory. Metrics are computed per endpoint before aggregation across starts. The central test uses fresh values in $[-3,3]$, and the outer test uses $[-10,-3]$ and $[3,10]$ with both signs for signed tasks. Positive self-power uses training inputs in $[0.02,3]$, central and development inputs drawn uniformly from $[0,3]$ (positive almost surely), and outer inputs in $[3,10]$. Each scalar evaluation has 2,048 samples per available outer side. Each arithmetic, sum or sorting test has 2,048 examples per split. The integer sequences use unseen indices 17--32 for early transfer and 33--64 for late transfer. The original one-shot study used 9--32 for early transfer. We preserve those scores separately because indices 9--16 are now used for feedback. Running sums in the primary comparison keep length six as magnitudes change. The supplement separately tests length 32. We report raw errors within tasks. For the aggregate ranking, we average raw outer RMSE over the sizes and starts of each task and then sum those ten task means. The supplementary plots also retain scaled errors and paired comparisons.
